\documentclass[12pt,a4paper]{article}
\usepackage[utf8]{inputenc}
\usepackage[spanish]{babel}
\usepackage{geometry}
\usepackage{listings}
\usepackage{xcolor}
\usepackage{parskip}
\usepackage{amsmath}
\usepackage{hyperref}

\geometry{margin=2.5cm}

\lstset{
  basicstyle=\ttfamily\small,
  keywordstyle=\color{blue}\bfseries,
  commentstyle=\color{gray},
  breaklines=true, frame=single, language=Java,
  numbers=left, numberstyle=\tiny\color{gray}
}

\title{\textbf{4-Puzzle Lineal --- BFS y DFS}\\
  \large Documentación Técnica}
\author{Inteligencia Artificial}
\date{}

\begin{document}
\maketitle
\tableofcontents
\newpage

\section{Descripción del Problema}
El \textbf{4-Puzzle Lineal} es un problema de búsqueda donde se tiene un arreglo de 4 elementos que deben quedar en orden ascendente $[1,2,3,4]$ usando únicamente tres operadores de intercambio.

\subsection{Operadores}
\begin{center}
\begin{tabular}{|l|l|l|}
\hline
\textbf{Operador} & \textbf{Acción} & \textbf{Posiciones} \\
\hline
Izquierdo & Intercambia pos 1$\leftrightarrow$2 & índices 0 y 1 \\
Central   & Intercambia pos 2$\leftrightarrow$3 & índices 1 y 2 \\
Derecho   & Intercambia pos 3$\leftrightarrow$4 & índices 2 y 3 \\
\hline
\end{tabular}
\end{center}

El espacio de búsqueda máximo es $4! = 24$ estados únicos.

\section{BFS --- Búsqueda por Amplitud}

Explora nivel a nivel usando una cola FIFO. Garantiza la ruta con el \textbf{mínimo número de movimientos}.

\begin{lstlisting}
int[] resolverBFS(int[] estadoInicial) {
    Queue<Nodo> cola = new LinkedList<>();
    Set<String> visitados = new HashSet<>();
    cola.add(new Nodo(estadoInicial, null));
    visitados.add(Arrays.toString(estadoInicial));

    while (!cola.isEmpty()) {
        Nodo actual = cola.poll();
        if (esMeta(actual.estado)) return actual;

        for (int[] hijo : generarHijos(actual.estado)) {
            String clave = Arrays.toString(hijo);
            if (!visitados.contains(clave)) {
                visitados.add(clave);
                cola.add(new Nodo(hijo, actual));
            }
        }
    }
    return null;
}
\end{lstlisting}

\section{DFS --- Búsqueda por Profundidad}

Explora profundamente usando una pila LIFO. No garantiza la solución óptima pero puede ser más rápido en algunos casos.

\begin{lstlisting}
int[] resolverDFS(int[] estadoInicial) {
    Stack<Nodo> pila = new Stack<>();
    Set<String> visitados = new HashSet<>();
    pila.push(new Nodo(estadoInicial, null));

    while (!pila.isEmpty()) {
        Nodo actual = pila.pop();
        String clave = Arrays.toString(actual.estado);
        if (visitados.contains(clave)) continue;
        visitados.add(clave);

        if (esMeta(actual.estado)) return actual;
        for (int[] hijo : generarHijos(actual.estado))
            pila.push(new Nodo(hijo, actual));
    }
    return null;
}
\end{lstlisting}

\section{Generación de Hijos}

Cada estado genera hasta 3 hijos aplicando los 3 operadores:

\begin{lstlisting}
List<int[]> generarHijos(int[] estado) {
    List<int[]> hijos = new ArrayList<>();
    for (int i = 0; i < 3; i++) { // 3 intercambios
        int[] hijo = estado.clone();
        int tmp = hijo[i];
        hijo[i] = hijo[i+1];
        hijo[i+1] = tmp;
        hijos.add(hijo);
    }
    return hijos;
}
\end{lstlisting}

\section{Visualización Animada}

Una vez obtenida la solución, se reconstruye el camino desde el nodo meta hasta el inicial siguiendo los punteros padre. La animación se implementa con \texttt{javax.swing.Timer}:

\begin{itemize}
  \item Cada 500ms se muestra el siguiente estado del camino.
  \item El tablero visual actualiza los colores de las celdas intercambiadas.
  \item Se muestra el número de paso actual sobre el total.
\end{itemize}

\section{Comparativa BFS vs DFS en este problema}

\begin{center}
\begin{tabular}{|l|c|c|}
\hline
\textbf{Característica} & \textbf{BFS} & \textbf{DFS} \\
\hline
Solución óptima & Sí & No siempre \\
Memoria usada & Mayor & Menor \\
Velocidad & Similar & Similar \\
Espacio explorado & Completo por nivel & Profundidad primero \\
\hline
\end{tabular}
\end{center}

\end{document}
